\documentclass[12pt,a4paper]{article}

% ---------------------------------------------------------
% Packages
% ---------------------------------------------------------
\usepackage[a4paper,
            left=1in,
            right=1in,
            top=0.8in,
            bottom=0.8in]{geometry}

\usepackage{newtxtext}
\usepackage{newtxmath}
\usepackage{pgfgantt}
\usepackage{graphicx}
\usepackage{amssymb}
\usepackage{array}
\usepackage{tabularx}
\usepackage{longtable}
\usepackage{booktabs}
\usepackage{multirow}
\usepackage{enumitem}
\usepackage{setspace}
\usepackage{fancyhdr}
\usepackage{titlesec}
\usepackage{caption}
\usepackage{float}
\usepackage{amssymb}
\usepackage{pifont}
\usepackage{ragged2e}
\usepackage{url}
\usepackage[hidelinks]{hyperref}

% ---------------------------------------------------------
% General Formatting
% ---------------------------------------------------------
\setlength{\parindent}{0pt}
\setlength{\parskip}{6pt}

\onehalfspacing

% Section formatting
\titleformat{\section}
  {\bfseries\large}
  {\thesection.}
  {0.5em}
  {}

\titlespacing*{\section}
  {0pt}
  {12pt}
  {6pt}

% ---------------------------------------------------------
% Header / Footer
% ---------------------------------------------------------
\pagestyle{fancy}
\fancyhf{}

\fancyfoot[R]{\thepage\;|\;P a g e}

\renewcommand{\headrulewidth}{0pt}
\renewcommand{\footrulewidth}{0pt}

% ---------------------------------------------------------
% Table formatting
% ---------------------------------------------------------
\renewcommand{\arraystretch}{1.25}

% ---------------------------------------------------------
% Custom commands
% ---------------------------------------------------------

% Green check mark similar to the original template
\newcommand{\greencheck}{%
    \colorbox{green!50}{\Large\checkmark}%
}

% Placeholder text
\newcommand{\placeholder}[1]{\textit{[#1]}}

% ---------------------------------------------------------
% Document
% ---------------------------------------------------------
\begin{document}

% =========================================================
% COVER PAGE
% =========================================================

\begin{titlepage}
\thispagestyle{empty}

\begin{center}

% ---------------------------------------------------------
% University Logo
% ---------------------------------------------------------
% Place the Varendra University logo in the same folder
% and name it "vu_logo.png".
%
% If you do not want to use a logo, comment the next line.
% ---------------------------------------------------------

\IfFileExists{vu_logo.png}{
    \includegraphics[width=1.0in]{vu_logo.png}
    \vspace{0.2in}
}{}

{\Large\bfseries Varendra University}

\vspace{0.08in}

{\large\bfseries Department of Computer Science and Engineering}

\vspace{0.55in}

{\large\bfseries Undergraduate Thesis/Project Proposal}

\vspace{0.55in}

{\large\bfseries Title of the Thesis/Project}

\vspace{0.08in}

{\itshape (Write your group proposal title)}

\vfill

{\bfseries Submitted by:}

\vspace{0.25in}

Name of the Student-1 (ID of the Student-1)\\
Name of the Student-2 (ID of the Student-2)\\
Name of the Student-3 (ID of the Student-3)

\vspace{0.25in}

%\begin{tabular}{rl}
\textbf{Program:} & B.Sc. in Computer Science and Engineering\\
\textbf{Semester:} & 7th\\
\textbf{Course Code:} & CSE 4100\\
\textbf{Course Title:} & Project or Thesis with Seminar part I
%\end{tabular}

\vfill

{\bfseries Supervisor:}

\vspace{0.08in}

[Full Name]\\
[Designation]\\
[Department Name]

\vfill

{\bfseries Submission Date:} [DD Month YYYY]

\end{center}

\end{titlepage}


% =========================================================
% PAGE 2
% =========================================================

\setcounter{page}{1}

\section{Introduction}

Introduce the topic, its background, and significance, highlighting its
relevance to current technologies \cite{james2021intro}. Ensure that the section adequately discusses the principles or computational frameworks of the proposed solution.

\vspace{0.1in}

\textit{(Approx. 200--300 words)}

\vspace{0.15in}


\section{Literature Review / Background Study}

Summarize 3--5 key existing works (research papers, systems, tools,
products). Explain their contributions and clearly identify gaps or
limitations that your work addresses.

Use proper citations --- IEEE style recommended like \cite{hasan2023thesis}, \cite{rahman2024ai}.


\section{Problem Statement}

Describe the specific problem or challenge your thesis/project will solve.
For a thesis, this should reflect a research gap \cite{uci2025dataset}. For projects, this may
involve real-world inefficiencies or limitations \cite{smith2023ml}.

\vspace{0.1in}

Use 1--2 focused paragraphs.


\section{Objectives}

List the intended goals of your work. These should be SMART --- Specific,
Measurable, Achievable, Relevant, and Time-bound.

\begin{itemize}[leftmargin=0.35in]
    \item Objective 1
    \item Objective 2
    \item Objective 3
    \item Objective 4
\end{itemize}


\section{Proposed Methodology / System Design}

Explain how you plan to achieve your objectives.

For a thesis, describe your research methods, models, and/or data.
For a project, define the system architecture, modules, tools, and
process flow.

\vspace{0.1in}

Include diagrams or flowcharts if applicable.


\section{Method Justification}

Explain and justify why the selected methodology/system design fits the
FYDP objectives/requirements.


\section{Expected Outcomes}

List or describe the main deliverables or findings you expect to achieve ---
e.g., working system, experimental results, research insights, prototype,
etc.


% =========================================================
% PAGE 3
% =========================================================

%\newpage

\section{Tools and Technologies to Be Used}

List the software, programming languages, libraries, hardware, or datasets
that will be used to implement or support your work.


\section{References}

\renewcommand{\refname}{}
\bibliographystyle{IEEEtran}
\bibliography{references}

\section{Work Division and Team Management}

Describe how tasks will be distributed among team members. Use a Table \ref{tab:work-division} to show task assignment. Clarify roles.

\begin{table}[H]
\centering
\caption{Work Division and Team Management}
\label{tab:work-division}
\begin{tabularx}{\textwidth}{|p{1.25in}|X|p{1.4in}|}
\hline
\textbf{Team Member} &
\textbf{Assigned Tasks} &
\textbf{Role} \\
\hline

Student-1 &
Literature review, methodology, implementation &
Member \\
\hline

Student-2 &
Data collection, system development, testing &
Member \\
\hline

Student-3 &
Documentation, analysis, presentation &
Member \\
\hline

\end{tabularx}
\end{table}


\section{Timeline}

Provide a time-based breakdown of your work (e.g., by weeks or months).
Use a table or Gantt chart.

\begin{figure}[H]
\centering
\resizebox{\textwidth}{!}{%
\begin{ganttchart}[
    hgrid,
    vgrid,
    x unit=1.0cm,
    y unit title=0.7cm,
    y unit chart=0.65cm,
    title height=1,
    bar height=0.5,
    bar label font=\small,
    title label font=\small\bfseries
]{1}{12}

\gantttitle{Thesis Research Timeline}{12} \\

\gantttitle{M1}{1}
\gantttitle{M2}{1}
\gantttitle{M3}{1}
\gantttitle{M4}{1}
\gantttitle{M5}{1}
\gantttitle{M6}{1}
\gantttitle{M7}{1}
\gantttitle{M8}{1}
\gantttitle{M9}{1}
\gantttitle{M10}{1}
\gantttitle{M11}{1}
\gantttitle{M12}{1} \\

\ganttbar{Literature Review}{1}{3} \\
\ganttbar{Research Problem and Objectives}{1}{2} \\
\ganttbar{Research Framework / Methodology}{3}{5} \\
\ganttbar{Data Collection and Preprocessing}{4}{6} \\
\ganttbar{Proposed Method Development}{5}{8} \\
\ganttbar{Experimental Investigation}{7}{9} \\
\ganttbar{Performance Evaluation}{8}{10} \\
\ganttbar{Results and Discussion}{9}{10} \\
\ganttbar{Thesis Writing}{9}{11} \\
\ganttbar{Revision and Finalization}{11}{12} \\
\ganttbar{Thesis Defense}{12}{12}

\end{ganttchart}%
}
\caption{Proposed 12-Month Thesis Timeline}
\label{fig:thesis_timeline}
\end{figure}

\begin{figure}[H]
\centering
\resizebox{\textwidth}{!}{%
\begin{ganttchart}[
    hgrid,
    vgrid,
    x unit=1.0cm,
    y unit title=0.7cm,
    y unit chart=0.65cm,
    title height=1,
    bar height=0.5,
    bar label font=\small,
    title label font=\small\bfseries
]{1}{12}

\gantttitle{Project Timeline}{12} \\

\gantttitle{M1}{1}
\gantttitle{M2}{1}
\gantttitle{M3}{1}
\gantttitle{M4}{1}
\gantttitle{M5}{1}
\gantttitle{M6}{1}
\gantttitle{M7}{1}
\gantttitle{M8}{1}
\gantttitle{M9}{1}
\gantttitle{M10}{1}
\gantttitle{M11}{1}
\gantttitle{M12}{1} \\

\ganttbar{Literature Review}{1}{2} \\

\ganttbar{Problem Identification}{1}{2} \\

\ganttbar{Requirement Analysis}{2}{3} \\

\ganttbar{System Design / Methodology}{3}{4} \\

\ganttbar{Data Collection and Preparation}{4}{5} \\

\ganttbar{System Implementation}{5}{8} \\

\ganttbar{Testing and Debugging}{7}{9} \\

\ganttbar{Performance Evaluation}{8}{10} \\

\ganttbar{Documentation}{9}{11} \\

\ganttbar{Final Review and Refinement}{11}{12} \\

\ganttbar{Final Presentation / Deployment}{12}{12}

\end{ganttchart}%
}
\caption{Proposed 12-Month Project Timeline}
\label{fig:project_timeline}
\end{figure}

\section{Cost-Benefit Analysis}

Estimate any costs for tools, software, hosting, datasets, or devices.
Mention if it is low-cost, cost-effective, or funded.

Assess the economic feasibility, compare different options, and select the
most cost-effective one. If no cost, state that clearly.


\section{Constraints and Limitations}

Mention only the constraints relevant to your work. These may include
economic, environmental, ethical, health and safety, social, political,
or sustainability-related factors that affect the design, development,
or research process.


\section{Societal Impact}

Briefly explain how your proposed work will benefit society. This may
include improvements in quality of life, accessibility, education,
public health, digital inclusion, environmental awareness, or other
social aspects.

Focus on real-world relevance and ethical contribution to the community
or specific user groups.


\section{Complex Engineering Problem Mapping}

Fill this table by reflecting how your proposal maps with the complex
engineering problem domains.

\begin{table}[H]
\centering
\caption{Mapping with complex problem solving.}
\label{tab:complex-problem}

\small
\renewcommand{\arraystretch}{1.3}

\begin{tabularx}{\textwidth}{|
>{\centering\arraybackslash}X|
>{\centering\arraybackslash}X|
>{\centering\arraybackslash}X|
>{\centering\arraybackslash}X|
>{\centering\arraybackslash}X|
>{\centering\arraybackslash}X|
>{\centering\arraybackslash}X|}
\hline

\textbf{WP1} &
\textbf{WP2} &
\textbf{WP3} &
\textbf{WP4} &
\textbf{WP5} &
\textbf{WP6} &
\textbf{WP7} \\
\hline

Depth of Knowledge &
Range of Conflicting Requirements &
Depth of Analysis &
Familiarity of Issues &
Extent of Applicable Codes &
Extent of Stakeholder Involvement &
Interdependence \\
\hline

{\Large\checkmark} &
{\Large\checkmark} &
{\Large\checkmark} &
{\Large\checkmark} &
&
&
\\
\hline

\end{tabularx}
\end{table}


\section{Complex Engineering Activity Mapping}

Justify which complex engineering activity level your proposal aligns with.

\begin{table}[H]
\centering
\caption{Mapping with complex engineering activities.}
\label{tab:complex-activity}

\small
\renewcommand{\arraystretch}{1.3}

\begin{tabularx}{\textwidth}{|
>{\centering\arraybackslash}X|
>{\centering\arraybackslash}X|
>{\centering\arraybackslash}X|
>{\centering\arraybackslash}X|
>{\centering\arraybackslash}X|}
\hline

\textbf{EA1} &
\textbf{EA2} &
\textbf{EA3} &
\textbf{EA4} &
\textbf{EA5} \\
\hline

Range of resources &
Level of Interaction &
Innovation &
Consequences for society and environment &
Familiarity \\
\hline

{\Large\checkmark} &
{\Large\checkmark} &
&
&
\\
\hline

\end{tabularx}
\end{table}


\section{Green Computing Considerations}

Mention how your work considers energy efficiency, environmental
sustainability, or digital carbon footprint reduction.

Even if indirectly related, reflect awareness of green ICT principles.


\section{Lifelong Learning and Personal Development}

State how this work will contribute to your knowledge, skills, and
long-term professional growth.

For example, learning new technologies, teamwork, research skills, etc.


\end{document}